\chapter{Gesture and Activity Recognition}

Everything we have built so far -- signals as vectors, the Fourier transform as a
projection onto sinusoids, spectrograms that localize frequency in time, filters
that sculpt a signal's spectrum, and the inertial and muscle sensors of the last
few chapters -- has been working toward a single concrete question. A band sits on
your forearm. A few sensors stream numbers many times a second. \emph{What is the
hand doing right now?} Is it making a fist, flicking left, tapping, rolling the
wrist, or just resting? And can we answer fast enough, and reliably enough, to
drive a musical instrument in real time?

That question is the subject of \textbf{gesture and activity recognition}. It is
the bridge between raw, wiggly sensor traces and meaningful, discrete labels:
\emph{fist}, \emph{wave}, \emph{walking}, \emph{typing}. The good news is that
almost none of it is new physics. It is a \emph{pipeline} built almost entirely
from tools you already own. We cut the endless stream into manageable
\textbf{windows}; from each window we compute a handful of \textbf{features}
(many of them just the time-domain and frequency-domain quantities from earlier
chapters); and we hand those features to a \textbf{classifier} that maps them to
a label. The art is in choosing the windows and features so that the geometry of
the problem -- where different gestures land in feature space -- becomes simple
enough that even a tiny classifier can draw the boundaries.

This chapter is clear about what this pipeline can and cannot do. A wrist IMU on
its own tells a flick from a fist from a roll. Add a force-myography (FMG) or
surface-EMG band on the forearm, and -- with a short per-finger calibration for
the person wearing it -- you can detect \emph{which finger moved}, because the
tendons and muscles that drive each finger leave a distinct mechanical and
electrical signature just below the band. What stays out of reach from the wrist
is the \emph{continuous joint angle} of each finger; that is genuinely
underdetermined and needs sensors on the fingers themselves. A recognizer tuned
to one person can also drop in accuracy on a stranger -- which is exactly why
\emph{calibration} sits at the center of our design rather than a one-size-fits-all
model. We will see \emph{why} each of these is true, and how to measure them so
you are never fooled by your own optimism.

\chapterline

\section{The recognition pipeline, end to end}
\label{ch38:pipeline}

Let us name the stages before we open any of them up. Sensor samples arrive as a
stream -- say a tri-axial accelerometer at $f_s = 100$ Hz, so every
$10$ ms we get a new triple $(a_x, a_y, a_z)$. We never classify a single sample;
one instant carries almost no information about a \emph{motion}, which by
definition unfolds over time. Instead:

\begin{enumerate}
  \item \textbf{Segmentation / windowing.} Group the stream into chunks of $N$
        samples. Each chunk -- a \emph{window} -- is the atomic unit we will
        label. (Section~\ref{ch38:windowing}.)
  \item \textbf{Feature extraction.} Reduce each window of (possibly many) raw
        samples to a short \emph{feature vector} $\vec{x} \in \R^{d}$ of
        descriptive numbers: how big the motion is, how fast it wiggles, how the
        axes move together, where its energy sits in frequency.
        (Sections~\ref{ch38:time-features}--\ref{ch38:freq-features}.)
  \item \textbf{Classification.} Map the feature vector to one of $K$ gesture
        labels, $\vec{x} \mapsto \hat{y} \in \set{1,\dots,K}$.
        (Section~\ref{ch38:classification}.)
  \item \textbf{Evaluation.} Measure how often $\hat{y}$ is right, on data the
        classifier has never seen, and ideally on \emph{people} it has never
        seen. (Section~\ref{ch38:evaluation}.)
\end{enumerate}

\begin{figure}[H]
\centering
\begin{tikzpicture}[
  font=\small,
  box/.style={draw, rounded corners, minimum height=1.0cm, minimum width=2.3cm,
              align=center, fill=blue!4},
  >=Stealth]
\node[box, fill=gray!8] (raw) {raw stream\\$(a_x,a_y,a_z)$};
\node[box, right=0.7cm of raw] (win) {window\\$N$ samples};
\node[box, right=0.7cm of win] (feat) {features\\$\vec{x}\in\R^d$};
\node[box, right=0.7cm of feat] (clf) {classifier};
\node[box, fill=green!6, right=0.7cm of clf] (lab) {label\\$\hat y$};
\draw[->] (raw) -- (win);
\draw[->] (win) -- (feat);
\draw[->] (feat) -- (clf);
\draw[->] (clf) -- (lab);
\end{tikzpicture}
\caption{The recognition pipeline. A continuous sensor stream is chopped into
windows; each window is reduced to a low-dimensional feature vector; a classifier
maps that vector to a discrete gesture label. Every box reuses tools from earlier
chapters.}
\label{ch38:fig-pipeline}
\end{figure}

\begin{intuitionbox}
The whole pipeline is one long act of \emph{dimensionality reduction with a
purpose}. A two-second accelerometer window at 100 Hz is a point in
$\R^{600}$ (200 samples $\times$ 3 axes) -- a space far too big and too noisy to
draw boundaries in. Feature extraction throws away the parts that do not matter
(exactly \emph{when} in the window a wiggle happened, the precise noise) and keeps
the parts that do (how strong, how fast, how coupled), landing us in a friendly
$\R^{10}$ or so. The classifier then just has to carve that small space into
gesture-shaped regions. Good features make those regions large and well
separated; bad features leave the classes tangled no matter how clever the
classifier.
\end{intuitionbox}

\begin{connectionbox}
Feature extraction is a \emph{lossy projection}, and you have met projections
throughout this book and the companion linear-algebra reader. The Fourier
transform projects onto sinusoids; PCA (Chapter on machine learning) projects
onto directions of maximum variance; here we project onto a hand-chosen set of
physically meaningful axes (energy, dominant frequency, inter-axis correlation).
The choice of \emph{what to project onto} is the entire game.
\end{connectionbox}

\begin{goalbox}
For the wrist controller this pipeline is split across two devices. A thin wrist
node samples the sensors, applies light filtering, compresses, and BLE-streams the
result; the paired phone does the heavy lifting -- it runs the guided per-finger
calibration, builds the user's templates, conditions the signals, and keeps the
streamed logs for deep offline analysis. For the low-latency live-play path a
lightweight template match runs on the node itself: every $\sim$0.1 s it emits a
verdict -- is a new gesture starting? which finger? -- and turns that into a MIDI
message (note on, pitch bend, control change) fast enough that a musician does not
feel the lag. That real-time constraint shapes every decision below: windows must
be short, features cheap, and the on-device match must fit in kilobytes. There is
no machine-learned model anywhere in the loop. The method is direct: calibrate
each user once, store their templates, and match the live signal against them.
\end{goalbox}

\section{Windowing and segmentation}
\label{ch38:windowing}

A window is just $N$ consecutive samples treated as one unit. Two numbers define
the windowing scheme: the \textbf{length} $N$ (or duration $T = N/f_s$) and the
\textbf{hop} $H$, the number of samples we advance before starting the next
window. The \emph{overlap} is $N - H$ samples, often quoted as a fraction
$1 - H/N$.

\begin{definition}[Sliding window]
\label{ch38:def-window}
Given a stream $x[0], x[1], \dots$, the $m$-th window of length $N$ with hop $H$
is the block
\[
  w_m = \big( x[mH],\, x[mH+1],\, \dots,\, x[mH + N - 1] \big), \qquad m = 0,1,2,\dots
\]
With $H = N$ the windows tile the stream without overlap; with $H < N$ they
overlap and a given sample contributes to several windows.
\end{definition}

\begin{figure}[H]
\centering
\begin{tikzpicture}[font=\small, >=Stealth]
% timeline
\draw[->] (0,0) -- (11.2,0) node[right] {sample index $n$};
\foreach \n in {0,1,...,10} \draw (\n,-0.06) -- (\n,0.06);
% windows length 4, hop 2
\draw[fill=blue!12, draw=blue!60!black] (0,0.25) rectangle (4,0.75)
  node[midway] {$w_0$};
\draw[fill=blue!8, draw=blue!60!black] (2,0.95) rectangle (6,1.45)
  node[midway] {$w_1$};
\draw[fill=blue!12, draw=blue!60!black] (4,1.65) rectangle (8,2.15)
  node[midway] {$w_2$};
\draw[fill=blue!8, draw=blue!60!black] (6,2.35) rectangle (10,2.85)
  node[midway] {$w_3$};
\node[align=left, anchor=west] at (0,3.4)
  {$N=4$ samples, hop $H=2 \Rightarrow$ 50\% overlap};
\end{tikzpicture}
\caption{Sliding windows of length $N=4$ with hop $H=2$ (50\% overlap). Each
window is classified independently; overlap means a gesture is ``looked at'' from
several alignments, which smooths the verdict and shortens latency.}
\label{ch38:fig-sliding}
\end{figure}

\subsection{How long should a window be?}

This is the first real design tension, and it is the same time--frequency
trade-off you met with the STFT (Chapter~10), wearing different clothes.

\begin{itemize}
  \item \textbf{Too short} and the window cannot contain a whole gesture. A flick
        of the wrist lasts maybe 300--500 ms; a 50 ms window catches only a
        sliver and cannot tell a flick from the start of a punch. Short windows
        also give the frequency-domain features almost no resolution (recall the
        DFT bin spacing is $f_s/N$).
  \item \textbf{Too long} and you blur distinct actions together, add latency
        (you cannot label a 2 s window until 2 s have elapsed), and risk a single
        window straddling two different gestures.
\end{itemize}

A useful rule of thumb: make the window roughly as long as the \emph{shortest
gesture you must distinguish}, then choose the hop to meet your latency budget.
For human activity and hand gestures, windows of $1$--$2.5$ s with $50\%$ overlap
are a common, defensible default; for the snappy, low-latency triggers a musician
needs, you push down toward $200$--$400$ ms and lean on detecting onsets rather
than steady states.

\begin{trapbox}
Beware analyzing a window long enough to contain a whole \emph{repetitive}
activity (walking, stirring) the same way you analyze a one-shot
\emph{gesture} (a flick). Repetitive activities have a clean periodicity that
frequency-domain features capture beautifully. One-shot gestures are
\emph{transients}: their information is in the \emph{shape} of a single burst, not
in any steady frequency. Using a long window and a centroid feature on a transient
throws away exactly what makes it recognizable. Match the analysis to the kind of
motion.
\end{trapbox}

\subsection{Event segmentation: finding where a gesture is}

Fixed sliding windows are simple and real-time friendly, but they are
\emph{dumb}: they chop on a clock, not on the action. Most windows of a wrist
stream contain nothing -- the hand is at rest -- and we waste effort classifying
emptiness, and worse, we risk slicing a gesture down the middle.

The fix is a cheap \textbf{onset detector} (or activity detector) that decides
\emph{when} something is happening, and only then fires the classifier. A robust
and almost free detector watches a short-time energy or jerk signal and triggers
when it crosses a threshold. Define the windowed energy of the (gravity-removed)
acceleration magnitude,
\begin{equation}
\label{ch38:energy}
  E[m] = \frac{1}{N}\sum_{n=0}^{N-1} \big\| \vec{a}[mH+n] - \bar{\vec{a}} \big\|^2 ,
\end{equation}
and call a window ``active'' when $E[m]$ exceeds a threshold $\theta$. To avoid
chattering at the threshold, use \emph{hysteresis}: a high threshold $\theta_{\text{on}}$
to start an event and a lower $\theta_{\text{off}}$ to end it. This is exactly the
Schmitt-trigger idea, and it is the difference between a controller that triggers
once per gesture and one that machine-guns a dozen spurious notes as your hand
crosses the threshold.

\begin{figure}[H]
\centering
\begin{tikzpicture}
\begin{axis}[
  width=12cm, height=4.6cm,
  xlabel={time (s)}, ylabel={short-time energy $E[m]$},
  xmin=0, xmax=6, ymin=0, ymax=1.25,
  axis lines=box, samples=200, domain=0:6,
  title={Energy-based event segmentation with hysteresis},
  title style={font=\small},
  legend style={at={(0.99,0.97)},anchor=north east, font=\scriptsize, draw=none, fill=none},
]
% energy: two bursts
\addplot[sigblue, thick, samples=200]
  {0.06 + 0.9*exp(-((x-1.6)^2)/0.10) + 0.75*exp(-((x-4.2)^2)/0.13)};
\addlegendentry{$E[m]$}
\addplot[trapred, dashed, thick, domain=0:6] {0.55};
\addlegendentry{$\theta_{\text{on}}$}
\addplot[quotientorange, dotted, thick, domain=0:6] {0.30};
\addlegendentry{$\theta_{\text{off}}$}
\end{axis}
\end{tikzpicture}
\caption{An onset detector with hysteresis. An event begins when energy rises
above $\theta_{\text{on}}$ and ends only when it falls below the lower
$\theta_{\text{off}}$. The gap between thresholds prevents a single gesture from
being split into many spurious triggers by noise near the boundary.}
\label{ch38:fig-onset}
\end{figure}

\begin{buildbox}
A minimal real-time segmenter on a microcontroller. Keep a running window energy;
emit a window for classification only on the rising edge through
\texttt{TH\_ON}, and refuse to re-trigger until energy has dropped below
\texttt{TH\_OFF}. This is a few lines and saves enormous CPU.
\begin{lstlisting}[language=C]
// streaming onset detector with hysteresis (one axis-magnitude sample at a time)
static int   armed = 1;          // ready to fire?
static float e = 0.0f;           // leaky energy estimate
const  float alpha = 0.05f;      // smoothing (approx 20-sample memory)
const  float TH_ON = 0.55f, TH_OFF = 0.30f;

int onset_step(float a_mag, float gravity) {
    float ac = a_mag - gravity;          // remove DC / gravity baseline
    e = (1.0f - alpha) * e + alpha * ac * ac;   // leaky short-time energy
    if (armed && e > TH_ON)  { armed = 0; return 1; }  // fire: classify now
    if (!armed && e < TH_OFF){ armed = 1; }            // re-arm
    return 0;
}
\end{lstlisting}
\end{buildbox}

\begin{historybox}
The threshold-with-hysteresis trick is the \emph{Schmitt trigger}, invented by
Otto Schmitt in 1934 while he was a graduate student studying nerve impulse
propagation in squid axons -- biology, not electronics. He noticed that a neuron
fires cleanly despite noisy, slowly-varying inputs precisely because it has two
different thresholds. The same idea now keeps your gesture controller from
stuttering. Good ideas in signal processing are frequently borrowed from the
nervous system, which has been doing real-time activity recognition for a very
long time.
\end{historybox}

\section{Time-domain features}
\label{ch38:time-features}

Now we reduce a window to numbers. Time-domain features are computed directly from
the samples without any transform -- they are cheap, robust, and surprisingly
powerful. Take a single-axis window $x[0],\dots,x[N-1]$ with mean
$\bar{x} = \frac1N\sum_n x[n]$.

\paragraph{Mean.} $\bar{x}$ itself. For an accelerometer it captures the gravity
component along that axis, hence \emph{orientation}: a flat wrist and a vertical
wrist have different per-axis means even when perfectly still. For a gyroscope or
high-pass-filtered signal the mean is near zero and carries little.

\paragraph{Variance / standard deviation.}
\begin{equation}
\label{ch38:variance}
  \sigma_x^2 = \frac{1}{N}\sum_{n=0}^{N-1}\big(x[n]-\bar{x}\big)^2 .
\end{equation}
The single most useful motion feature. It measures \emph{how much} the signal
moves regardless of orientation, separating ``still'' from ``active'' and gentle
from vigorous. Standing and sitting both have low variance; walking and waving
have high variance.

\paragraph{RMS (root mean square).}
\begin{equation}
\label{ch38:rms}
  \mathrm{RMS}(x) = \sqrt{\frac{1}{N}\sum_{n=0}^{N-1} x[n]^2}.
\end{equation}
The signal's effective magnitude including its DC level; note
$\mathrm{RMS}^2 = \bar{x}^2 + \sigma_x^2$, so RMS blends ``how tilted'' with
``how much it moves.'' For zero-mean signals RMS equals the standard deviation and
is a clean energy proxy.

\paragraph{Zero-crossing rate (ZCR).} The number of times the signal crosses its
mean per window,
\begin{equation}
\label{ch38:zcr}
  \mathrm{ZCR} = \frac{1}{N-1}\sum_{n=1}^{N-1}
     \tfrac12\,\big| \operatorname{sgn}(x[n]-\bar x) - \operatorname{sgn}(x[n-1]-\bar x)\big|.
\end{equation}
A poor-man's frequency estimate computed without a single multiply: a fast wiggle
crosses its mean often, a slow one rarely. We met ZCR in the audio-features
chapter; here it cheaply distinguishes a brisk tremor from a slow roll.

\paragraph{Min, max, range, and peak count.} Order statistics of the window. The
\emph{range} $\max - \min$ captures the amplitude of a one-shot transient that
variance smears out; the number of peaks above a threshold counts repetitions
(steps, taps).

\paragraph{Correlation between axes.} Real gestures move several axes \emph{in a
coordinated way}, and that coordination is a fingerprint. The Pearson correlation
between axes $x$ and $y$ over the window,
\begin{equation}
\label{ch38:corr}
  \rho_{xy} = \frac{\sum_{n}(x[n]-\bar x)(y[n]-\bar y)}
                   {\sqrt{\sum_n (x[n]-\bar x)^2}\,\sqrt{\sum_n (y[n]-\bar y)^2}}
            = \frac{\inner{x-\bar x}{\,y-\bar y}}{\norm{x-\bar x}\,\norm{y-\bar y}},
\end{equation}
is exactly the cosine of the angle between the two mean-centered window vectors --
the inner-product geometry from Chapter 3. $\rho_{xy} \approx +1$ means the axes
rise and fall together (a diagonal motion), $-1$ means they oppose, $0$ means they
are unrelated. A wrist \emph{roll} couples two axes strongly; a straight
\emph{flick} drives mostly one. Correlation is what lets a classifier tell the
\emph{direction} of a motion, not just its size.

\begin{whybox}
Why is inter-axis correlation so informative when variance per axis is not?
Because variance is \emph{rotation-blind in a misleading way}: two completely
different motions can deposit the same total variance, just distributed
differently across axes and time. Correlation looks at the \emph{joint} behavior
-- the \emph{shape} of the cloud of $(x[n],y[n])$ points. A motion confined to a
tilted line in the $xy$-plane has $|\rho| \approx 1$; an isotropic shake fills a
disk and has $\rho \approx 0$. That shape is a geometric signature of the gesture
that survives changes in how hard you perform it.
\end{whybox}

\begin{figure}[H]
\centering
\begin{tikzpicture}
\begin{axis}[
  width=12cm, height=5.2cm,
  xlabel={time (s)},
  xmin=0, xmax=2, ymin=-1.6, ymax=1.6,
  axis lines=box, samples=150, domain=0:2,
  title={Two gestures, same per-axis variance, different correlation},
  title style={font=\small},
  legend style={at={(0.99,0.03)},anchor=south east, font=\scriptsize, draw=none, fill=none},
]
% coupled (roll): x and y in phase -> rho near +1
\addplot[sigblue, thick] {sin(deg(2*pi*2*x))*exp(-0.3*x)};
\addlegendentry{axis $x$ (roll)}
\addplot[sigblue, thick, dashed] {0.9*sin(deg(2*pi*2*x)+10)*exp(-0.3*x)};
\addlegendentry{axis $y$ (roll): $\rho\approx +0.97$}
% anti / independent
\addplot[specred, thick] {0.5+sin(deg(2*pi*2*x)-180)*exp(-0.3*x)};
\addlegendentry{axis $x$ (flick)}
\addplot[specred, thick, dashed] {0.5+0.2*sin(deg(2*pi*5*x))*exp(-0.3*x)};
\addlegendentry{axis $y$ (flick): $\rho\approx 0$}
\end{axis}
\end{tikzpicture}
\caption{Per-axis variance alone cannot separate these two gestures, but their
inter-axis correlation does: the roll drives $x$ and $y$ together
($\rho \approx +1$), while the flick drives $x$ alone ($\rho \approx 0$).
Correlation encodes the \emph{direction} of motion.}
\label{ch38:fig-corr}
\end{figure}

\begin{connectionbox}
The list above -- variance, RMS, correlation -- is the elementary statistics of
the companion probability book applied to a window. Variance is
$\E[(X-\E X)^2]$ estimated over the window; correlation is the normalized
covariance $\operatorname{Cov}(X,Y)/(\sigma_X\sigma_Y)$. Treating a window as a
small sample from a random process is not a metaphor; it is the literal bridge
between Chapter 20 (random signals, noise) and feature engineering. A feature is
an \emph{estimator} of some property of the underlying motion, and like all
estimators it has bias and variance of its own.
\end{connectionbox}

\section{Frequency-domain features}
\label{ch38:freq-features}

For \emph{periodic} or \emph{textured} motions -- walking, running, scrubbing, a
sustained tremor -- the most discriminating information lives in the spectrum.
Take the DFT of the window, $X[k] = \sum_{n=0}^{N-1} x[n]\,e^{-\jj 2\pi kn/N}$, and
form the magnitude spectrum $|X[k]|$ over bins $k = 0,\dots,N/2$, with bin $k$
corresponding to frequency $f_k = k f_s / N$.

\paragraph{Band energies.} Partition the frequency axis into a few bands and sum
the power in each:
\begin{equation}
\label{ch38:band}
  E_{\text{band}} = \sum_{k \in \text{band}} |X[k]|^2 .
\end{equation}
Human locomotion lives below $\sim$3 Hz (step cadence) with harmonics up to
$\sim$10 Hz; voluntary hand gestures are slower still; physiological and pathological
tremor sits in characteristic bands (rest tremor $\approx 4$--6 Hz, essential
tremor $\approx 6$--12 Hz). A handful of band energies is often enough to separate
walking from running from idle.

\paragraph{Dominant frequency.} The bin with the most power,
$f^\star = f_{\,\arg\max_k |X[k]|}$, e.g. the step rate while walking. Cheap and
very informative for periodic activities; nearly meaningless for one-shot
transients (whose energy is spread, not peaked).

\paragraph{Spectral centroid.} The power-weighted average frequency -- the
spectrum's ``center of mass'',
\begin{equation}
\label{ch38:centroid}
  f_{\text{c}} = \frac{\sum_k f_k\,|X[k]|^2}{\sum_k |X[k]|^2}.
\end{equation}
We met the centroid as ``brightness'' in the audio-features chapter; for motion it
says whether the energy is concentrated low (slow, sweeping) or high (jittery,
fast). A gentle wave has a low centroid; a nervous shake has a high one.

\paragraph{Spectral entropy.} How \emph{spread out} the spectrum is. Normalize
the power spectrum to a distribution $p_k = |X[k]|^2 / \sum_j |X[j]|^2$ and take
\begin{equation}
\label{ch38:entropy}
  S = -\sum_k p_k \log p_k .
\end{equation}
Low entropy means the energy is concentrated at one frequency -- a clean periodic
motion like steady walking. High entropy means the energy is smeared across many
frequencies -- broadband, irregular, noise-like motion. Entropy beautifully
separates ``rhythmic'' from ``erratic'' even when the dominant frequency is the
same.

\begin{figure}[H]
\centering
\begin{tikzpicture}
\begin{axis}[
  width=6.0cm, height=4.4cm,
  xlabel={freq (Hz)}, ylabel={$|X[k]|^2$},
  xmin=0, xmax=10, ymin=0, ymax=1.1,
  axis lines=box, ybar, bar width=4pt,
  title={Walking: low entropy},
  title style={font=\small},
  xtick={0,2,...,10}, ytick=\empty,
]
\addplot[fill=fillsig, draw=sigblue] coordinates {
  (0,0.03)(1,0.08)(2,1.0)(3,0.10)(4,0.35)(5,0.06)(6,0.18)(7,0.04)(8,0.05)(9,0.03)(10,0.02)};
\end{axis}
\end{tikzpicture}
\hspace{0.4cm}
\begin{tikzpicture}
\begin{axis}[
  width=6.0cm, height=4.4cm,
  xlabel={freq (Hz)}, ylabel={$|X[k]|^2$},
  xmin=0, xmax=10, ymin=0, ymax=1.1,
  axis lines=box, ybar, bar width=4pt,
  title={Fidgeting: high entropy},
  title style={font=\small},
  xtick={0,2,...,10}, ytick=\empty,
]
\addplot[fill=specred!40, draw=specred] coordinates {
  (0,0.30)(1,0.45)(2,0.40)(3,0.55)(4,0.48)(5,0.50)(6,0.42)(7,0.52)(8,0.38)(9,0.44)(10,0.35)};
\end{axis}
\end{tikzpicture}
\caption{Two power spectra with similar total energy. Walking concentrates power
at the step frequency and its harmonics (low spectral entropy); irregular
fidgeting spreads power across the band (high entropy). The dominant-frequency
feature would call the right plot's peak meaningless, but entropy cleanly flags it
as broadband.}
\label{ch38:fig-spectra}
\end{figure}

\begin{trapbox}
Do not compute frequency features on a window without first removing the mean (DC)
and ideally applying a taper (Hann window). The huge DC bin from gravity will
dominate $\sum |X[k]|^2$ and drag the centroid toward 0 Hz, and spectral leakage
from the rectangular cutoff (Chapter 10) will smear sharp peaks into neighbors,
corrupting band energies. Mean-remove, taper, then transform.
\end{trapbox}

\section{A worked example: four wrist gestures}
\label{ch38:worked}

Let us make this concrete with small, hand-computable numbers. Suppose we capture
short windows from four gestures on a wrist IMU and (after preprocessing) read off
just three features per window: the variance of the $x$-axis acceleration
$\sigma_x^2$, the dominant frequency $f^\star$, and the $xy$ correlation
$\rho_{xy}$. Idealized values:

\begin{center}
\begin{tabular}{lccc}
\toprule
Gesture & $\sigma_x^2$ (g$^2$) & $f^\star$ (Hz) & $\rho_{xy}$ \\
\midrule
Rest        & 0.002 & --   & $\approx 0$ \\
Flick left  & 0.45  & 2.5  & $+0.1$ \\
Wrist roll  & 0.30  & 2.0  & $+0.9$ \\
Tremor      & 0.05  & 8.0  & $+0.2$ \\
\bottomrule
\end{tabular}
\end{center}

Look at how cleanly these separate, \emph{one feature at a time}: variance alone
peels Rest off the bottom; the dominant frequency isolates Tremor at the top;
correlation distinguishes the otherwise-similar Flick and Roll. No single feature
does the whole job, but together they place the four gestures at four corners of a
small feature cube, and \emph{any} reasonable classifier can separate corners.

\begin{example}[Computing features from raw samples]
\label{ch38:ex-compute}
Take a tiny window of $x$-axis acceleration (gravity already removed), $N = 8$:
\[
  x = (0.0,\; 0.4,\; 0.6,\; 0.2,\; -0.3,\; -0.5,\; -0.2,\; 0.1)\ \text{g}.
\]
Mean: $\bar x = \tfrac{1}{8}(0.3) = 0.0375 \approx 0$. Variance (about the mean),
\[
  \sigma_x^2 = \tfrac18\sum (x[n]-\bar x)^2 \approx \tfrac18(0.939) \approx 0.117\ \text{g}^2,
\]
so $\mathrm{RMS} \approx \sigma_x \approx 0.34$ g. Zero crossings of $x-\bar x$:
the sign pattern is $-,+,+,+,-,-,-,+$, which flips three times (between samples
0--1, 3--4 and 6--7), so $\mathrm{ZCR} = 3/(N-1) = 3/7 \approx 0.43$ crossings per
sample. At $f_s = 100$ Hz that suggests a fundamental of a few Hz -- a quick
sanity check against whatever the DFT reports. Every one of these is a few
additions; this is why time-domain features are the backbone of an embedded
recognizer.
\end{example}

\begin{labbox}
Record yourself (or just simulate) doing each of the four gestures ten times.
Compute the three features per window and scatter-plot variance against
correlation, coloring by gesture. You should \emph{see} the four clusters with
your own eyes -- and where two clusters overlap is exactly where your classifier
will make its mistakes. This plot, made before you train anything, is the single
most valuable diagnostic in the whole project: if the clusters are not visually
separable in your features, no classifier will rescue you, and you should go back
and find better features.
\end{labbox}

\section{Classification}
\label{ch38:classification}

Once each window is a feature vector $\vec{x} \in \R^d$, classification is the act
of carving $\R^d$ into $K$ labeled regions. You do not need a deep network. For
the small, well-separated feature spaces of gesture recognition, simple
classifiers are fast, interpretable, and frequently \emph{better} (less prone to
overfitting on the modest data you can realistically collect from a few people).

\paragraph{$k$-nearest neighbors (k-NN).} Store all labeled training feature
vectors. To classify a new $\vec{x}$, find the $k$ stored points nearest to it
(Euclidean distance, on \emph{normalized} features) and take a majority vote. It
is a one-line idea -- ``it is whatever its neighbors are'' -- with no training at
all, and it makes the geometry of Section~\ref{ch38:worked} literal: a gesture is
recognized by landing near other examples of the same gesture. Its cost is memory
and per-query distance computations, which can sting on a microcontroller.

\paragraph{Decision trees and random forests.} A tree asks a sequence of
threshold questions -- ``is $\sigma_x^2 > 0.1$? if so, is $f^\star > 5$?'' --
each splitting feature space with an axis-aligned cut, until a leaf assigns a
label. This is \emph{exactly} the one-feature-at-a-time reasoning we did by eye in
Section~\ref{ch38:worked}, automated and made optimal. A \emph{random forest}
(many trees voting) is one of the most reliable off-the-shelf classifiers for
tabular features like ours, and small trees run in microseconds on a
microcontroller.

\paragraph{Linear and logistic classifiers.} Draw a hyperplane
$\vec{w}\cdot\vec{x} + b = 0$ between classes; the sign of $\vec{w}\cdot\vec{x}+b$
is the decision. When classes are linearly separable in feature space (often the
case after good feature engineering) this is the cheapest possible classifier --
one dot product -- and its weights $\vec{w}$ are interpretable as ``how much each
feature matters.''

\begin{figure}[H]
\centering
\begin{tikzpicture}
\begin{axis}[
  width=10.5cm, height=6.0cm,
  xlabel={variance $\sigma_x^2$ (normalized)}, ylabel={correlation $\rho_{xy}$},
  xmin=-0.1, xmax=1.1, ymin=-0.3, ymax=1.1,
  axis lines=box,
  title={Feature space with a calibrated decision boundary},
  title style={font=\small},
  legend style={at={(0.99,0.5)},anchor=east, font=\scriptsize, draw=none, fill=none},
]
% rest cluster (low var, low corr)
\addplot[only marks, sourceblue, mark=*] coordinates {
  (0.02,0.05)(0.04,-0.02)(0.01,0.10)(0.05,0.03)(0.03,0.07)};
\addlegendentry{rest}
% flick (high var, low corr)
\addplot[only marks, specred, mark=triangle*, mark size=2.5] coordinates {
  (0.85,0.12)(0.92,0.05)(0.80,0.18)(0.95,0.02)(0.88,0.10)};
\addlegendentry{flick}
% roll (high var, high corr)
\addplot[only marks, targetgreen, mark=square*] coordinates {
  (0.70,0.88)(0.78,0.92)(0.65,0.85)(0.82,0.90)(0.72,0.95)};
\addlegendentry{roll}
% a decision boundary illustration
\draw[structurepurple, thick, dashed] (axis cs:0.35,-0.3) -- (axis cs:0.35,1.1);
\draw[structurepurple, thick, dashed] (axis cs:0.35,0.5) -- (axis cs:1.1,0.5);
\end{axis}
\end{tikzpicture}
\caption{Three gestures in a two-feature space. Good features make the classes
fall into separated clusters; a tiny decision tree -- ``high variance? then high
correlation?'' -- draws the two axis-aligned cuts (dashed) that separate them. The
classifier is doing exactly the by-eye reasoning of Section~\ref{ch38:worked}.}
\label{ch38:fig-decision}
\end{figure}

\begin{focusbox}
\textbf{Always normalize your features before classifying.} A k-NN or linear model
measures distances in feature space; if variance ranges over $[0, 0.5]$ and
dominant frequency over $[0, 50]$, the frequency axis silently dominates every
distance and the variance feature is ignored. Standardize each feature to zero
mean and unit standard deviation (using statistics computed \emph{on the training
set only}), then store those constants and apply the same transform on the device.
This single step often matters more than the choice of classifier.
\end{focusbox}

\begin{happeningbox}
What is really happening when a classifier fits to data? It is drawing the
\emph{boundaries} between the clusters you already created with your features. The
classifier adds almost nothing if the clusters are clean -- a child could draw the
lines in Figure~\ref{ch38:fig-decision}. Conversely, no classifier can separate
clusters that overlap. This is why experienced practitioners spend their time on
windowing and features, not on swapping classifiers: \textbf{the features decide
whether the problem is solvable; the classifier merely solves the solvable
problem.} For our controller we lean on the simplest end of this spectrum -- store
the user's calibration examples and label a new window by the nearest stored
template -- precisely because clean features make the boundaries trivial and a
direct match needs no training procedure at all.
\end{happeningbox}

\section{Template matching and dynamic time warping}
\label{ch38:dtw}

The feature-vector approach summarizes a window into static numbers, discarding
the order in which things happened. For some gestures the \emph{trajectory} --
the precise temporal shape -- is the point: a checkmark motion, a circle, a
signature. For these, \textbf{template matching} keeps the whole trajectory and
compares an incoming gesture to stored exemplars.

The naive comparison is the Euclidean distance between two sequences sample by
sample. It fails badly because two performances of the same gesture are never the
same length or speed: you do the circle a little faster this time, or pause in the
middle. Aligning sample $n$ of one to sample $n$ of the other compares the top of
one circle to the side of another and reports a huge, meaningless distance.

\textbf{Dynamic time warping (DTW)} fixes this by finding the best
\emph{nonlinear} alignment between two sequences before measuring their distance.
Given sequences $a$ (length $M$) and $b$ (length $N$), it builds a warping path
through the $M\times N$ grid of local distances $d(i,j) = (a_i - b_j)^2$ that
stretches and compresses time to match the shapes. The optimal-path cost is found
by dynamic programming:
\begin{equation}
\label{ch38:dtw-rec}
  D(i,j) = d(i,j) + \min\big\{\, D(i-1,j),\; D(i,j-1),\; D(i-1,j-1)\,\big\},
\end{equation}
with $D(0,0)=d(0,0)$ along the boundaries, and the DTW distance is $D(M-1,N-1)$.
The three predecessors correspond to ``advance in $a$ only'' (compress $b$),
``advance in $b$ only'' (stretch $b$), and ``advance both'' (match). The cheapest
path through the grid is the best way to line up the two gestures despite
differences in tempo.

\begin{figure}[H]
\centering
\begin{tikzpicture}[scale=0.55, font=\scriptsize, >=Stealth]
% grid
\draw[gray!50] (0,0) grid (7,5);
\node[below] at (3.5,-0.4) {sequence $a$ (time)};
\node[rotate=90, above] at (-0.5,2.5) {sequence $b$ (time)};
% a monotone warping path
\draw[structurepurple, ultra thick, ->]
  (0.5,0.5) -- (1.5,0.5) -- (1.5,1.5) -- (2.5,1.5) -- (3.5,2.5)
  -- (4.5,3.5) -- (4.5,4.5) -- (6.5,4.5);
% endpoints
\fill[targetgreen] (0.5,0.5) circle (4pt);
\fill[trapred] (6.5,4.5) circle (4pt);
\node[right] at (7.2,4.5) {end};
\node[right] at (0.7,0.15) {start};
\end{tikzpicture}
\caption{A dynamic time warping path through the cost grid. The path moves
monotonically from corner to corner; horizontal/vertical steps stretch or
compress one sequence in time to match the other, while diagonal steps match
samples directly. DTW measures shape similarity while tolerating differences in
speed.}
\label{ch38:fig-dtw}
\end{figure}

\begin{intuitionbox}
Picture two people humming the same tune, one slightly slower with a pause in the
middle. Lay one performance along the bottom of a grid and the other up the side.
A diagonal line would assume they are perfectly in sync; DTW instead finds a
\emph{wiggly monotone staircase} that pairs each note of one with the matching
note of the other, pausing (going straight along a row) exactly where one singer
paused. The total mismatch \emph{along that best staircase} is the DTW distance.
It answers ``are these the same shape, played at possibly different speeds?''
\end{intuitionbox}

\begin{whybox}
Why does the simple recurrence~\eqref{ch38:dtw-rec} find the \emph{globally}
optimal alignment? Because the best path to cell $(i,j)$ must pass through one of
its three valid predecessors, and the best path to \emph{each} of those was
already computed and stored. This is the optimal-substructure property that makes
dynamic programming work: a globally optimal alignment is built from optimal
sub-alignments. The cost is $O(MN)$ time and memory -- fine for a stored library
of short templates, but it grows with the number and length of templates, which
is the practical limit on DTW.
\end{whybox}

\begin{trapbox}
DTW is seductive but does not scale. Matching every incoming gesture against a
large library is $O(MN)$ \emph{per template}, and DTW distances are not a metric
you can index easily. For a handful of distinctive whole-trajectory gestures (a
few templates per class) it is excellent and needs no training. For many classes,
many users, or high throughput, the window-feature-classifier path of
Sections~\ref{ch38:time-features}--\ref{ch38:classification} is usually the better
engineering choice. Use a constraint band (Sakoe--Chiba) to cap how far the path
may stray from the diagonal: it speeds DTW up and, by forbidding absurd
alignments, often improves accuracy too.
\end{trapbox}

\begin{applicationbox}
DTW was born in 1970s speech recognition (Sakoe and Chiba, 1978), where it aligned
a spoken word against templates despite people talking at different speeds. The
same algorithm now powers gesture recognition on smartwatches, signature
verification, gait analysis, and bioinformatics sequence alignment. Whenever you
need ``same shape, different speed,'' DTW is the classical answer -- a single
elegant idea that crossed from speech into motion sensing wholesale.
\end{applicationbox}

\section{Evaluating a recognizer}
\label{ch38:evaluation}

A recognizer that has not been carefully evaluated is just a hope. This section is
about not fooling yourself -- the most important and most neglected skill in the
whole pipeline.

\subsection{Accuracy and its lies}

\textbf{Accuracy} is the fraction of windows labeled correctly. It is the obvious
metric and it is frequently misleading. If 95\% of your windows are ``rest'' (as
they will be -- hands are still most of the time), a classifier that
\emph{always} says ``rest'' scores 95\% accuracy while detecting exactly zero
gestures. With imbalanced classes, accuracy rewards the trivial.

\subsection{The confusion matrix}

The full summary is the \textbf{confusion matrix} $C$, where $C_{ij}$ counts
windows whose true label is $i$ but were predicted $j$. The diagonal is correct
predictions; every off-diagonal entry is a specific, named mistake. The confusion
matrix tells you not just \emph{how often} you are wrong but \emph{how} -- and
``flick mistaken for roll'' is an actionable clue (add a correlation feature)
while a bare accuracy number is not.

\begin{figure}[H]
\centering
\begin{tikzpicture}[font=\small]
\matrix (m) [matrix of nodes,
  nodes={minimum size=1.05cm, anchor=center},
  column sep=-\pgflinewidth, row sep=-\pgflinewidth,
  nodes in empty cells]
{
  |[fill=green!30]| 48 & |[fill=red!10]| 1  & |[fill=red!10]| 1  \\
  |[fill=red!10]| 2    & |[fill=green!30]| 44 & |[fill=red!25]| 4 \\
  |[fill=red!10]| 1    & |[fill=red!25]| 6  & |[fill=green!30]| 43 \\
};
\draw (m-1-1.north west) rectangle (m-3-3.south east);
\foreach \r in {1,2} \draw (m-\r-1.south west) -- (m-\r-3.south east);
\foreach \c in {1,2} \draw (m-1-\c.north east) -- (m-3-\c.south east);
% labels
\node[above=2pt] at (m-1-1.north) {pred R};
\node[above=2pt] at (m-1-2.north) {pred F};
\node[above=2pt] at (m-1-3.north) {pred Roll};
\node[left=2pt] at (m-1-1.west) {true R};
\node[left=2pt] at (m-2-1.west) {true F};
\node[left=2pt] at (m-3-1.west) {true Roll};
\end{tikzpicture}
\caption{A confusion matrix for three gestures (Rest, Flick, Roll). The strong
diagonal is good; the highlighted off-diagonal cells -- Flick predicted as Roll
(4) and Roll predicted as Flick (6) -- reveal that these two are the classifier's
main confusion, exactly as Figure~\ref{ch38:fig-corr} would predict if the
correlation feature were weak or noisy.}
\label{ch38:fig-confusion}
\end{figure}

From the confusion matrix come the per-class metrics that survive class imbalance.
For a class $c$: \textbf{precision} $= \frac{\text{TP}}{\text{TP}+\text{FP}}$
(of windows you \emph{called} $c$, how many were?), \textbf{recall}
$= \frac{\text{TP}}{\text{TP}+\text{FN}}$ (of windows that \emph{were} $c$, how
many did you catch?), and their harmonic mean the \textbf{F1 score}
$= 2\,\frac{\text{precision}\cdot\text{recall}}{\text{precision}+\text{recall}}$.
For a controller, precision and recall trade off in a way you can \emph{feel}: low
precision means spurious notes you did not intend; low recall means gestures you
performed that the instrument missed. Which is worse depends on the instrument.

\subsection{The cardinal sin: leaky evaluation}

You must measure performance on data the classifier never saw in training. That
much is well known. What is less appreciated -- and what silently inflates most
reported gesture-recognition numbers -- is \emph{how} you split.

\begin{trapbox}
\textbf{Never split overlapping windows randomly into train and test.} With 50\%
overlap, adjacent windows share half their samples; a random split puts a window
in training and its near-twin in test, so the classifier is graded on data it has
effectively memorized. Accuracy looks spectacular and means nothing. Split by
\emph{recording session} or by \emph{time block}, never by individual window.
\end{trapbox}

\subsection{Cross-subject generalization: the number that hurts}

Here is the truth the glossy demos omit. Fit and test on the \emph{same person}
and you will see very high accuracy -- often 95\%+ -- because the templates match
\emph{that person's} particular way of moving. Now test on a \emph{different}
person, using someone else's templates, and accuracy commonly falls hard, into the
70s or worse. The way to measure this gap is \textbf{leave-one-subject-out
cross-validation (LOSO)}: fit on everyone except one person, test on that held-out
person, and repeat for each person. The LOSO number is what a stranger would
experience if you handed them a band with \emph{someone else's} templates -- and it
is precisely the gap our per-user calibration step is designed to close. We do not
ship a stranger's boundaries; we record the new user's own finger movements at
enrollment, so the templates always belong to the hand wearing the band.

\begin{figure}[H]
\centering
\begin{tikzpicture}
\begin{axis}[
  width=11cm, height=5.0cm,
  ybar, bar width=22pt,
  symbolic x coords={within-subject, mixed (leaky), cross-subject},
  xtick=data, x tick label style={font=\scriptsize, align=center},
  ymin=0, ymax=115, ylabel={accuracy (\%)},
  axis lines=box,
  title={The same recognizer, three evaluation protocols},
  title style={font=\small},
  nodes near coords, nodes near coords style={font=\scriptsize},
  enlarge x limits=0.25,
]
\addplot[fill=fillsig, draw=sigblue] coordinates {
  (within-subject,96) (mixed (leaky),93) (cross-subject,74)};
\end{axis}
\end{tikzpicture}
\caption{Illustrative (and typical) accuracies for one recognizer under three
protocols. Within-subject and the leaky mixed split look great; the
leave-one-subject-out (cross-subject) number is the one a new user feels with
someone else's templates. Always report LOSO -- the others flatter your model.}
\label{ch38:fig-loso}
\end{figure}

\begin{whybox}
Why does accuracy drop so much across people? Because the feature distributions
\emph{shift}. People differ in limb length, where the band sits, how they wear it
(rotation!), baseline tremor, strength, and personal motion style. A boundary
tuned to your variance and correlation values lands in the wrong place for
someone whose ``vigorous flick'' is another person's ``gentle one.'' This is
\emph{distribution shift}, and it is the central difficulty of cross-subject
wearable recognition. The cleanest remedy is the one we build the whole product
around: a brief per-user calibration so the templates and thresholds belong to
\emph{this} wearer, not to an average of strangers. Supporting tactics help too --
orientation-invariant features (use acceleration \emph{magnitude} and inter-axis
correlation rather than raw axes), per-user normalization, and re-running the
calibration whenever the band is re-seated. With calibration, the cross-subject
gap stops being your problem: every user effectively runs a within-subject
recognizer.
\end{whybox}

\begin{connectionbox}
Cross-subject drop is a concrete instance of the train/test mismatch the machine
learning chapter formalizes as the difference between empirical risk (error on
your sample) and true risk (error on the real population). LOSO is your best cheap
estimator of true risk for the population ``people who are not in your dataset.''
It usually has high variance with few subjects, so report the mean \emph{and}
spread across the held-out folds, not a single rosy figure.
\end{connectionbox}

\section{What the wrist can and cannot sense}
\label{ch38:limits}

We have built a pipeline that recognizes a vocabulary of distinct
wrist-and-forearm motions. Now let us draw the line precisely between what is
within reach and what is not -- because the line is not where casual intuition
puts it.

\begin{focusbox}
\textbf{Which finger moved is detectable from the wrist; each finger's continuous
joint angle is not.} The accelerometer and gyroscope sense the wrist as a rigid
segment, so on their own they give you whole-hand and wrist gestures (fist vs.\
open, flick, twist, tap-induced shock, double-tap). But the muscles and tendons
that drive the fingers run \emph{through the forearm, directly under a band}. Tap
your index versus your ring finger and the tendon bulge and electrical activity
just below the band differ in a repeatable way. A force-myography or surface-EMG
sensor picks that up, and with a short per-finger calibration the band detects
\emph{which finger moved}. What remains genuinely out of reach from the wrist is
the \emph{continuous flexion angle} of each joint: that information is simply not
present in wrist-frame bulk motion, and recovering it needs sensors on the fingers.
\end{focusbox}

To detect \emph{which finger} acted, add a forearm sensing modality alongside the
IMU:
\begin{itemize}
  \item \textbf{Force myography (FMG)} reads the mechanical bulging of forearm
        muscles with a ring of pressure sensors (FSRs). It is cheap, robust to
        electrical noise and sweat, and after a short per-finger enrollment it
        separates finger and grip patterns well (Figure~\ref{ch38:fig-fmg}). This
        is the principle behind WristFlex, which distinguished finger pinches from a
        wrist FMG band.
  \item \textbf{Surface EMG} on the forearm senses the electrical activity of the
        finger-driving muscles, the basis of myoelectric control (Myo, MyoWare).
        It catches \emph{intended} movement slightly before motion appears, but is
        more sensitive to electrode shift and sweat than FMG, so it benefits most
        from per-user calibration and re-seating checks.
  \item \textbf{Continuous finger \emph{pose}} -- actual joint angles for a piano-like
        ten-voice instrument -- requires sensors \emph{on the fingers} (rings with
        IMUs, stretch sensors, instrumented gloves). Recent finger-worn work such
        as ``Learning on the Rings'' tracks continuous flexion this way. There is no
        way around the physics: to measure a joint you must observe something that
        the joint actually moves. We treat this as an optional upgrade path, not a
        requirement for the instrument we are building.
\end{itemize}

\begin{figure}[H]
\centering
\begin{tikzpicture}
\begin{axis}[
  width=12cm, height=5.0cm,
  xlabel={FSR channel around the forearm (1--8)}, ylabel={pressure (normalized)},
  xmin=0.5, xmax=8.5, ymin=0, ymax=1.15,
  axis lines=box, xtick={1,2,...,8},
  title={Per-finger FMG signatures from an 8-channel forearm band},
  title style={font=\small},
  legend style={at={(0.99,0.97)},anchor=north east, font=\scriptsize, draw=none, fill=none},
]
% index finger: bulge concentrated on channels 2-3
\addplot[sigblue, thick, mark=*, mark size=1.4] coordinates {
  (1,0.18)(2,0.85)(3,0.62)(4,0.22)(5,0.10)(6,0.08)(7,0.12)(8,0.15)};
\addlegendentry{index tap}
% ring finger: bulge concentrated on channels 5-6
\addplot[specred, thick, dashed, mark=square*, mark size=1.4] coordinates {
  (1,0.12)(2,0.16)(3,0.20)(4,0.40)(5,0.80)(6,0.70)(7,0.30)(8,0.14)};
\addlegendentry{ring tap}
% rest: flat baseline
\addplot[targetgreen, thick, dotted] coordinates {
  (1,0.10)(2,0.11)(3,0.10)(4,0.12)(5,0.10)(6,0.11)(7,0.10)(8,0.11)};
\addlegendentry{rest}
\end{axis}
\end{tikzpicture}
\caption{Force-myography snapshots across an eight-channel forearm band. Tapping
the index finger bulges the tendons under channels~2--3; tapping the ring finger
bulges channels~5--6; at rest the profile is flat. After a short per-finger
enrollment, matching the live channel vector to these stored signatures tells
\emph{which finger} moved -- something a wrist IMU alone cannot see. The signatures
are person-specific, which is exactly why calibration is per user.}
\label{ch38:fig-fmg}
\end{figure}

\begin{goalbox}
This sets the design of our controller. From the wrist IMU we get a reliable
\emph{discrete gesture vocabulary} (taps, flicks, twists, holds) plus continuous
wrist \emph{orientation}, a superb expressive pitch-bend or modulation axis. Add a
forearm FMG/FSR array (and optionally surface EMG) and, after the per-finger
calibration, the band reports \emph{which finger} acted -- enough to play notes and
chords by finger. The edge over single-sensor lab demos comes from \emph{fusing}
modalities: a high-rate IMU for gross motion and orientation, the FMG/FSR array for
finger identity, and a separate high-rate accelerometer used as a contact /
bio-acoustic pickup that hears the sharp transient of a fingertip tap (the ViBand
idea). No one of these is enough alone; together, with per-user calibration, they
cover finger \emph{events} cleanly. The one thing we do not promise is ten
continuously variable finger voices like a piano -- that needs finger-worn sensors.
For everything a player actually does on this instrument -- triggering notes and
chords with finger movements, bending and modulating with the wrist -- the wrist
band delivers.
\end{goalbox}

\begin{historybox}
Wearable activity recognition matured in the 2000s as MEMS accelerometers became
cheap and tiny. Bao and Intille's 2004 study, recognizing everyday activities from
multiple body-worn accelerometers, is a landmark that established the
window-feature-classifier template still used today. The cross-subject gap
they and successors documented has not been ``solved'' so much as managed; it
remains the reason a fitness tracker counts \emph{your} steps well after a few days
and a stranger's less so. The thread that leads to our wrist instrument runs
through specific systems: WristFlex (finger pinch detection from a wrist FMG band),
the Myo and MyoWare surface-EMG modules, ViBand (UIST 2016, a high-rate
wrist accelerometer used as a bio-acoustic pickup for tap and contact events), a
PLOS One 2020 study fusing IMU with acoustic sensing, an MDPI Sensors 2023 study
counting gestures from a wrist IMU, and ``Learning on the Rings'' for finger-worn
continuous tracking. Our edge over any single one of these is to fuse modalities --
IMU, FMG/FSR array, and a bio-acoustic accelerometer -- and to calibrate per finger
per user, rather than chase one sensor with a heavier model.
\end{historybox}

\section{Putting it together: a reference recognizer}
\label{ch38:reference}

Here is the whole pipeline in a few lines of NumPy -- windowing, a compact feature
vector mixing time- and frequency-domain quantities, and a nearest-template match
against the per-user enrollment. It is deliberately small and correct; it is the
seed of what we will port to the microcontroller in the embedded chapter.

\begin{lstlisting}[language=Python]
import numpy as np

def features(win, fs):
    """win: (N,3) tri-axial window (gravity removed). Returns a feature vector."""
    feats = []
    for ax in range(3):                       # per-axis time-domain features
        x = win[:, ax]
        feats += [x.mean(), x.var(), np.sqrt((x**2).mean())]   # mean, var, RMS
        zc = np.mean(np.abs(np.diff(np.sign(x - x.mean())))) / 2
        feats.append(zc)                                       # zero-crossing rate
    # inter-axis correlations (the coupling / direction signature)
    c = np.corrcoef(win.T)                    # 3x3 correlation matrix
    feats += [c[0,1], c[0,2], c[1,2]]
    # frequency-domain on the acceleration magnitude
    mag = np.linalg.norm(win, axis=1)
    mag = (mag - mag.mean()) * np.hanning(len(mag))   # detrend + taper
    P = np.abs(np.fft.rfft(mag))**2
    f = np.fft.rfftfreq(len(mag), 1/fs)
    P[0] = 0                                   # kill residual DC
    centroid = np.sum(f * P) / (np.sum(P) + 1e-12)
    dom = f[np.argmax(P)]
    p = P / (P.sum() + 1e-12)
    entropy = -np.sum(p * np.log(p + 1e-12))
    feats += [centroid, dom, entropy]
    return np.array(feats)

def windows(stream, N, H):
    """Yield (start, window) for each sliding window."""
    for m in range((len(stream) - N) // H + 1):
        s = m * H
        yield s, stream[s:s+N]

# usage sketch (direct template matching -- no learned model):
#   T = {g: features(cal_win, fs=100) for g, cal_win in enrollment}  # per-finger templates
#   x = (features(w, fs=100) - mu) / sigma     # normalize with enrollment mu, sigma
#   y_hat = min(T, key=lambda g: np.linalg.norm(x - (T[g]-mu)/sigma))  # nearest template
\end{lstlisting}

\begin{aestheticbox}
Step back and notice what we have. Twenty-odd numbers per window -- means,
variances, RMS, zero-crossings, three correlations, a centroid, a dominant
frequency, an entropy -- and almost every one is something you derived from first
principles in an earlier chapter. The inner product became correlation; the DFT
became band energy and centroid; the STFT's time--frequency trade-off became the
window-length decision; the probability book's variance and expectation became
features. Gesture recognition is not a separate subject bolted on at the end. It
is the \emph{payoff}: the same handful of geometric ideas, pointed at your moving
hand.
\end{aestheticbox}

\section{Summary}
\label{ch38:summary}

\begin{itemize}
  \item Recognition is a pipeline: \emph{window} the stream, extract a low-dim
        \emph{feature vector}, \emph{classify} it, \emph{evaluate} without fooling
        yourself.
  \item Window length trades latency and gesture-containment against frequency
        resolution -- the STFT trade-off in a new guise. An onset detector with
        hysteresis segments events cheaply and avoids spurious triggers.
  \item Time-domain features (mean, variance, RMS, ZCR, range, inter-axis
        correlation) are cheap and powerful; correlation, the cosine between
        mean-centered axis vectors, encodes \emph{direction} of motion.
  \item Frequency-domain features (band energies, dominant frequency, spectral
        centroid, spectral entropy) shine for periodic/textured motion; always
        remove DC and taper first.
  \item Simple classifiers (k-NN, decision trees/forests, linear) suffice when
        features separate the classes; features matter more than the classifier.
        Always normalize features with train-set statistics.
  \item DTW matches whole trajectories despite speed differences via dynamic
        programming; great for a few template gestures, poor at scale.
  \item Evaluate with confusion matrices and per-class precision/recall/F1, never
        bare accuracy on imbalanced classes; never split overlapping windows
        randomly; report leave-one-subject-out accuracy as the number that counts.
  \item A wrist IMU yields whole-hand gestures plus orientation. Adding a forearm
        FMG/EMG band and a short per-finger calibration detects \emph{which finger}
        moved. What the wrist cannot give is each finger's continuous joint angle;
        that needs sensors on the fingers.
\end{itemize}

\section{Exercises}
\label{ch38:exercises}

\begin{exercise}
A sensor streams at $f_s = 100$ Hz. You choose windows of $N = 256$ samples with
$50\%$ overlap. (a) What is the window duration in seconds? (b) What is the hop in
samples and in seconds? (c) How many windows per second does the classifier
process? (d) What is the DFT bin spacing $f_s/N$ in Hz, and what is the lowest
nonzero frequency you can resolve?
\end{exercise}

\begin{exercise}
For the 8-sample window in Example~\ref{ch38:ex-compute}, verify the variance
$\sigma_x^2 \approx 0.117$ by computing $\sum (x[n]-\bar x)^2$ by hand. Then
compute the range ($\max-\min$) and explain in one sentence which feature --
variance or range -- you would trust more to flag a sharp, one-shot tap, and why.
\end{exercise}

\begin{exercise}
Two accelerometer axes over a window are $x = (1,2,3,2,1)$ and $y = (2,4,6,4,2)$.
(a) Compute the Pearson correlation $\rho_{xy}$ from \eqref{ch38:corr}. (b) Without
recomputing, state $\rho$ for $y' = (-2,-4,-6,-4,-2)$ and explain geometrically
why scaling and sign flips affect (or do not affect) the correlation.
\end{exercise}

\begin{exercise}
You build an onset detector with a single threshold $\theta$ (no hysteresis). The
energy hovers right around $\theta$ during a slow gesture, dithering above and
below it due to noise. Describe the failure you will observe at the controller's
output, and explain quantitatively how splitting $\theta$ into
$\theta_{\text{on}} > \theta_{\text{off}}$ fixes it. What is the cost of making the
gap too wide?
\end{exercise}

\begin{exercise}
A dataset is 92\% ``rest'' and 8\% ``gesture'' windows. A lazy classifier always
predicts ``rest.'' (a) What is its accuracy? (b) What are its precision and recall
\emph{for the gesture class}? (c) What is its F1 for the gesture class? Use this to
argue in two sentences why F1 (or a confusion matrix) is the right metric here.
\end{exercise}

\begin{exercise}
Compute a small DTW distance. Let $a = (0, 2, 0)$ and $b = (0, 0, 2, 2, 0)$ with
local cost $d(i,j) = |a_i - b_j|$. Fill in the $3\times 5$ cumulative cost grid
using \eqref{ch38:dtw-rec} and report $D$ at the top-right corner. Sketch (or
describe) the optimal warping path and say which samples of $b$ get matched to the
peak of $a$.
\end{exercise}

\begin{exercise}
You report 96\% accuracy by randomly splitting 50\%-overlapping windows into 80\%
train / 20\% test. A colleague says the number is meaningless. (a) Explain the
specific leak. (b) Describe two concrete changes to the split that would make the
estimate trustworthy. (c) Predict qualitatively how the reported accuracy will change
after each fix, and why.
\end{exercise}

\begin{exercise}
Design a five-feature vector to distinguish four wrist gestures: \emph{rest},
\emph{flick}, \emph{wrist roll}, and \emph{double-tap}. For each feature name it,
say whether it is time- or frequency-domain, and state in one phrase which pair of
gestures it primarily helps separate. Then identify which two gestures you expect
to be hardest to tell apart and what extra feature or sensor would help.
\end{exercise}

\begin{exercise}
(Cross-subject reasoning.) You train on five people and test LOSO, getting per-fold
accuracies $\{0.81, 0.62, 0.88, 0.55, 0.79\}$. (a) Report the mean and the spread.
(b) The fold at 0.55 is one specific held-out person; give two physical reasons
that one user might be much harder than the others. (c) Propose one change to the
feature set that targets the most likely cause, and explain why it should help.
\end{exercise}

\begin{exercise}
(Real-time budget.) Your microcontroller must, every hop, compute the feature
vector and run a small decision tree, all within the hop interval. With
$N = 200$, which features in the reference code of Section~\ref{ch38:reference} are
$O(N)$ and which are $O(N\log N)$? If you must cut cost, which single feature would
you drop first and which would you never drop, and why?
\end{exercise}

\begin{exercise}
(Sensing-limits check.) A teammate proposes that, using only the wrist IMU, the
controller will continuously output each of the five fingers' joint angles to
drive five independent synth voices. State plainly whether this is achievable with
a wrist IMU alone, explain the physical reason in one or two sentences, and
describe the minimal sensing change that would make a \emph{discrete} per-finger
scheme (which finger tapped) feasible. Then say which sensing change is required to
recover the \emph{continuous} joint angles, and why it cannot live at the wrist.
\end{exercise}
